% INSERT BEFORE: online_appendix.tex "\section{D. Institutional pilot and empirical design}\label{oa-d}"
% Source: the C.10 contract in the numerical worktree, paper/online_appendix.md (commit 9e7ab4d).
% The contract's \input of the table is replaced by a pointer: the exhibit is Table 4 of the paper.
% \NUM{name} is filled from the quantity registry; \INP{case.key} from numerics/manifests/c10_classes.json.
\subsection{C.10. Quality classes, general priors, and benchmark
windows}\label{oa-c-classes}

\textbf{Inputs.} This exercise uses the general form of Proposition 2
of the paper. Its high-cost conditions are (A1), and its trade
conditions are (A2). The exact declarations appear in the notes to
Table 4 of the paper and in the run manifest. Cases A and B each change
three benchmark inputs: the high cost, the trade cost, and the weak
strength. The other two strengths stay fixed. Case C uses uniform
classes above the reserve, with the same means as the point masses.
Case D replaces those classes by point masses and retains every other
input.

The posterior bounds are

\[
m(\pi)=\frac{\pi}{\pi+(1-\pi)e^{2/b}},\qquad
M(\pi)=\frac{\pi}{\pi+(1-\pi)e^{-2/b}},\qquad
\kappa(\pi)=\min\{m(\pi),1-M(\pi)\}.
\]

Class payoffs are the exact integrals of the auction polynomials. The
high-class win probability is
\(W_H(r)=\mathbb E_H[\min\{\theta/r,1\}]\), since both supports exceed
the reserve. At an interior threshold,

\[
\tau=\frac{c_H-g_L}{g_H-g_L},\qquad
x^*=\frac b2\left[\log\frac{\tau}{1-\tau}-\log\frac{\pi}{1-\pi}\right].
\]

Conditional entry follows the Laplace tails at \(x^*-1\) and
\(x^*+1\). Total entry is \(\pi e_H+(1-\pi)e_L\), and high-class
ownership is \(\pi e_HW_H\). No-trade rows have no flow threshold. A
threshold above \(M(\pi)\) is unattainable; equality retains the
positive mass on the upper plateau.

\textbf{Conditions and status.} Each case records every parameter
restriction, both class supports, their order, and equal admissibility.
The record also gives all four margins in (A1)--(A2), the collapse
margin, and the trade bound at the collapse strength. The sufficient
ownership margin is \((1+\rho)W_H(r_1)-2\rho W_H(r_0)\). The exact
ownership margin is \(e_H(r_1)W_H(r_1)-\rho W_H(r_0)\). Gross profit at
the lowest posterior is a diagnostic; zero cost gives preparation
without an extra strict condition. Cases A--C satisfy the complete
proposition, including part (iii) and the sufficient ownership
condition.

Case D fails the upper high-cost condition at the strong strength. The
record retains this negative margin. Its trade bounds still imply unique
no trade at the weak strength and unique full orders at both higher
strengths. High-cost preparation is infeasible at both higher strengths,
so total entry stays at its floor. Analytical status comes from these
global bounds; an order mesh supplies only a numerical diagnostic.
Table 4 of the paper reports the outcomes and the margins of all four
cases.

\textbf{Benchmark intervals.} Three different prior intervals answer
three different questions. The trade-cost window is nonempty for priors
in approximately
\((\NUM{base_prior_window_lower},\NUM{base_prior_window_upper})\). With
both benchmark costs fixed, all hypotheses for parts (i)--(ii) hold on
\((\NUM{base_prior_prop2_lower},\NUM{base_prior_prop2_upper})\). With the
collapse strength fixed as well, parts (i)--(iii) require
\((\NUM{base_prior_prop2_all_lower},\NUM{base_prior_prop2_all_upper})\).
These intervals are open; the calculations use the unrounded endpoints.
The broader first interval does not establish the proposition at the
fixed benchmark costs.

At the equal prior, the trade-cost window is \((1/60,\,1/[12(1+e)])\),
approximately
\((\NUM{base_trade_cost_lower},\NUM{base_trade_cost_upper})\). The fixed
cost is \(k=\INP{benchmark.k}\). The spread ratio is
\NUM{base_spread_ratio}, against the sufficient bound
\NUM{base_spread_ratio_bound}. Across the four noise laws with atomic
costs, strong entry ranges from \NUM{four_noise_entry_min} to
\NUM{four_noise_entry_max}. The hyperbolic secant and Student values
retain their status as diagnostics at fixed full orders. This range
establishes no additional equilibrium claim.

\textbf{Validation and outputs.} Exact decimal strings enter fifty-digit
arithmetic directly. Independent integration over realized bids
validates target receipts, challenger profits, and class win
probabilities. Independent density roots and full-line integrals
validate the threshold, conditional entry, mean posterior, and mean
price. Held-schedule deviations use both signs and both declared order
grids, with the required tighter quadrature pass. The analytical Laplace
tails have zero truncation error. All C.0 tolerances stay unchanged; a
breach leaves the result open and stops the producer.

The outputs and the run manifest are:

\begin{Shaded}
\begin{Highlighting}[]
\NormalTok{numerics/class\_exhibit.csv}
\NormalTok{numerics/class\_conditions.csv}
\NormalTok{numerics/class\_bounds.csv}
\NormalTok{numerics/class\_deviations.csv}
\NormalTok{numerics/class\_diagnostics.json}
\NormalTok{numerics/manifests/c10\_classes.json}
\end{Highlighting}
\end{Shaded}

The manifest records complete inputs, methods, tolerances, software
versions, source hashes, output hashes, and acceptance. The registry and
table consume only these validated CSV files. The full gate repeats the
calculation and compares every saved row. Feed: Table 4 and Sections
4.2 and 5.2 of the paper.

